% p233_1; source PDF page 233; manually verified editable table.
\begin{tikzpicture}[x=1cm,y=-1cm,inner sep=0pt,outer sep=0pt]
\path[use as bounding box] (0,0) rectangle (16.000,18.650);
\draw[line width=.4pt] (0.000,0.000) rectangle (2.000,0.850);
\node[align=center,text width=1.740cm,font={\linespread{1}\fontsize{10}{13.50}\selectfont\bfseries}] at (1.000,0.425) {行为类型};
\draw[line width=.4pt] (2.000,0.000) rectangle (4.200,0.850);
\node[align=center,text width=1.940cm,font={\linespread{1}\fontsize{10}{13.50}\selectfont\bfseries}] at (3.100,0.425) {主观};
\draw[line width=.4pt] (4.200,0.000) rectangle (7.100,0.850);
\node[align=center,text width=2.640cm,font={\linespread{1}\fontsize{10}{13.50}\selectfont\bfseries}] at (5.650,0.425) {结论};
\draw[line width=.4pt] (7.100,0.000) rectangle (9.900,0.850);
\node[align=center,text width=2.540cm,font={\linespread{1}\fontsize{10}{13.50}\selectfont\bfseries}] at (8.500,0.425) {罪数};
\draw[line width=.4pt] (9.900,0.000) rectangle (16.000,0.850);
\node[align=center,text width=5.840cm,font={\linespread{1}\fontsize{10}{13.50}\selectfont\bfseries}] at (12.950,0.425) {范例};
\draw[line width=.4pt] (0.000,0.850) rectangle (2.000,8.850);
\node[align=left,text width=1.740cm,font={\linespread{1}\fontsize{10}{13.50}\selectfont}] at (1.000,4.850) {非法拘禁的实行行为（整体是一个行为）};
\draw[line width=.4pt] (2.000,0.850) rectangle (4.200,3.350);
\node[align=left,text width=1.940cm,font={\linespread{1}\fontsize{10}{13.50}\selectfont}] at (3.100,2.100) {对重伤、死亡结果持过失心理};
\draw[line width=.4pt] (4.200,0.850) rectangle (7.100,3.350);
\node[align=left,text width=2.640cm,font={\linespread{1}\fontsize{10}{13.50}\selectfont}] at (5.650,2.100) {非法拘禁罪致人重伤、死亡（结果加重犯）};
\draw[line width=.4pt] (7.100,0.850) rectangle (9.900,3.350);
\node[align=left,text width=2.540cm,font={\linespread{1}\fontsize{10}{13.50}\selectfont}] at (8.500,2.100) {一个行为，最终只定一个罪};
\draw[line width=.4pt] (9.900,0.850) rectangle (16.000,3.350);
\node[align=left,text width=5.840cm,font={\linespread{1}\fontsize{10}{13.50}\selectfont}] at (12.950,2.100) {甲用绳子捆绑了乙，次日忽然发现乙窒息死亡。甲构成非法拘禁罪致人死亡};
\draw[line width=.4pt] (2.000,3.350) rectangle (4.200,8.850);
\node[align=left,text width=1.940cm,font={\linespread{1}\fontsize{10}{13.50}\selectfont}] at (3.100,6.100) {对重伤、死亡结果持故意心理（例如放任）};
\draw[line width=.4pt] (4.200,3.350) rectangle (7.100,8.850);
\node[align=left,text width=2.640cm,font={\linespread{1}\fontsize{10}{13.50}\selectfont}] at (5.650,6.100) {非法拘禁罪与故意伤害罪（重伤）、故意杀人罪想象竞合，择一重罪论处};
\draw[line width=.4pt] (7.100,3.350) rectangle (9.900,8.850);
\node[align=left,text width=2.540cm,font={\linespread{1}\fontsize{10}{13.50}\selectfont}] at (8.500,6.100) {一个行为，最终只定一个罪};
\draw[line width=.4pt] (9.900,3.350) rectangle (16.000,8.850);
\node[align=left,text width=5.840cm,font={\linespread{1}\fontsize{10}{13.50}\selectfont}] at (12.950,6.100) {甲为了非法拘禁乙，大力反扭乙的胳膊，致其胳膊折断，导致重伤。此时伤害行为是为了拘禁乙，也是非法拘禁罪的实行行为。甲的一个伤害行为同时触犯非法拘禁罪和故意伤害罪（重伤），想象竞合，择一重罪论处，定故意伤害罪（重伤）};
\draw[line width=.4pt] (0.000,8.850) rectangle (2.000,18.650);
\node[align=left,text width=1.740cm,font={\linespread{1}\fontsize{10}{13.50}\selectfont}] at (1.000,13.750) {非法拘禁的实行行为之外的暴力行为（前后有两个行为）};
\draw[line width=.4pt] (2.000,8.850) rectangle (4.200,13.650);
\node[align=left,text width=1.940cm,font={\linespread{1}\fontsize{10}{13.50}\selectfont}] at (3.100,11.250) {过失致人重伤、死亡};
\draw[line width=.4pt] (4.200,8.850) rectangle (7.100,13.650);
\node[align=left,text width=2.640cm,font={\linespread{1}\fontsize{10}{13.50}\selectfont}] at (5.650,11.250) {（法律拟制）故意伤害罪（重伤）、故意杀人罪};
\draw[line width=.4pt] (7.100,8.850) rectangle (9.900,13.650);
\node[align=left,text width=2.540cm,font={\linespread{1}\fontsize{10}{13.50}\selectfont}] at (8.500,11.250) {两个行为。非法拘禁行为即使已经构成非法拘禁罪，最终也只定故意伤害罪（重伤）、故意杀人罪};
\draw[line width=.4pt] (9.900,8.850) rectangle (16.000,13.650);
\node[align=left,text width=5.840cm,font={\linespread{1}\fontsize{10}{13.50}\selectfont}] at (12.950,11.250) {甲对乙已经构成非法拘禁罪后，乙辱骂甲，甲想教训乙，扇耳光，碰巧打到乙的太阳穴，导致乙死亡。对甲拟制为故意杀人罪};
\draw[line width=.4pt] (2.000,13.650) rectangle (4.200,18.650);
\node[align=left,text width=1.940cm,font={\linespread{1}\fontsize{10}{13.50}\selectfont}] at (3.100,16.150) {故意重伤、杀人};
\draw[line width=.4pt] (4.200,13.650) rectangle (7.100,18.650);
\node[align=left,text width=2.640cm,font={\linespread{1}\fontsize{10}{13.50}\selectfont}] at (5.650,16.150) {（正常处理）故意伤害罪（重伤）、故意杀人罪};
\draw[line width=.4pt] (7.100,13.650) rectangle (9.900,18.650);
\node[align=left,text width=2.540cm,font={\linespread{1}\fontsize{10}{13.50}\selectfont}] at (8.500,16.150) {两个行为。非法拘禁行为若构成非法拘禁罪，则两个行为应数罪并罚};
\draw[line width=.4pt] (9.900,13.650) rectangle (16.000,18.650);
\node[align=left,text width=5.840cm,font={\linespread{1}\fontsize{10}{13.50}\selectfont}] at (12.950,16.150) {甲对乙已经构成非法拘禁罪后，乙辱骂甲，甲生气，便杀了乙。甲构成非法拘禁罪和故意杀人罪。多数说认为，两罪应并罚。少数说认为，故意杀人罪吸收非法拘禁罪（2021年、2022年主观题）};
\end{tikzpicture}
